\scanpage{80}
On a ainsi \deleted{des \unsure{isomorphismes}}\inserted{des \uline{lois covariantes}}
pour la cohomologie. C'est très explicite pour
$f$ lisse. Pour $f$ général, pour l'expliciter,
on \deleted{\unsure{en}} \unsure{se servira parfois de}

2$^\circ$)\quad \uline{L'isom de transitivité}
\[
\mathbb{R}^!(gf)\simeq \mathbb{R}^!f\,\mathbb{R}^!g
\]
(Pour \unsure{dém.}, cf plus bas).

3$^\circ$)\quad Revenons alors à 1$^\circ$) lorsque $X,Y$
sont lisses sur $Z$, soit $F^\bullet\in\operatorname{Ob}\mathrm D^+(Z)$,
et prenons \deleted{$\unsure{G^\bullet}$}
\[
G^\bullet=\mathbb{L}g^*(F^\bullet)
\]
d'où \inserted{si $g$ est partout de dim rel. $d$,}
\[
\deleted{\unsure{g^!}}G^\bullet\simeq g^!(F^\bullet)\deleted{[\unsure{-2d|2d}]}
 \otimes T_{Y/Z}[-2d]
\]
\[
f^!(G^\bullet)=
 \underbrace{[f^!g^!(F^\bullet)]}_{\substack{
  \simeq\,h^!(F^\bullet)\\
  \Vert\\\mathbb{L}h^*(F^\bullet)\deleted{[2d']}\inserted{\otimes T_{X/Z}}[2d']}}
 \otimes T_{Y/Z}^{-1}[-2d]
\]
\margin{si \deleted{\unsure{h}} \unsure{partout}
de dim relative $d'$}
\[
f^!(G^\bullet)\simeq \mathbb{L}h^*(F^\bullet)\otimes T_{X/Y}
 [\unsure{+|\,}2\unsure{d_{X/Y}}]
\]
où
\[
T_{X/Y}\simeq T_{X/Z}\otimes T_{Y/Z}^{-1}
\]
\[
d_{X/Y}=\dim_{X/Z}-\dim_{Y/Z}
 =\text{dim \uline{\unsure{relative absolue}}}.
\]
\mspage{79}
